% ---------------------------------------------------------------------------
% Figure: scoring object detection as a two-step process
%   (1) assign predicted boxes to truth boxes, then
%   (2) read the assignment off as a binary-classification confusion matrix.
%
% Photograph: "cat-and-dog.jpg"
%   Source: https://www.publicdomainpictures.net/en/view-image.php?image=24076&picture=cat-and-dog
%   License: CC0 1.0 (Public Domain Dedication). No attribution required; credited here as a courtesy.
% ---------------------------------------------------------------------------
\documentclass[tikz,border=4pt]{standalone}
\usepackage[T1]{fontenc}
\usepackage{lmodern}
\usepackage{graphicx}
\usepackage{xcolor}
\usepackage{amsmath}
\usepackage{amssymb}
\usepackage{tikz}
\usetikzlibrary{arrows.meta,positioning,calc}

\definecolor{kwgreen}{HTML}{3EAE2B}
\definecolor{kwblue}{HTML}{0068C7}
\definecolor{kworange}{HTML}{EF7724}
\definecolor{kwred}{HTML}{F42836}
\definecolor{kwdarkgreen}{HTML}{2E5524}
\definecolor{kwdarkblue}{HTML}{003765}

% Identity colors for the truth / prediction boxes. 
% OUTSIDE the green/orange/red/blue family used for the confusion outcomes, so a
% "truth box" is never mistaken for a TP and a "prediction box" never for a TN.
% These are placeholders -- tweak them freely.
\definecolor{truthcol}{HTML}{5E3A9E}   % truth boxes     
\definecolor{predcol}{HTML}{A65A7A}    % prediction boxes 

% Slate gray: 5A6473



\newcommand{\TP}{\textcolor{kwgreen}{\texttt{\textbf{TP}}}}
\newcommand{\FP}{\textcolor{kwred}{\texttt{\textbf{FP}}}}
\newcommand{\FN}{\textcolor{kworange}{\texttt{\textbf{FN}}}}
\newcommand{\TN}{\textcolor{kwblue}{\texttt{\textbf{TN}}}}

% Displayed size of each photo copy in Panel 1 (the photo is 1280x1920, i.e. 2:3 portrait).
\def\Wd{1.80}   % width  in cm
\def\Hd{2.70}   % height in cm  (= 1.80 * 1920/1280)

\begin{document}
\begin{tikzpicture}[
    >=Latex,
    font=\sffamily,
    panel/.style={draw=black!18, rounded corners=2pt, line width=0.7pt, fill=white},
    ptitle/.style={font=\bfseries\normalsize, text=kwdarkblue, anchor=north west},
    body/.style={font=\scriptsize, text=black},
    note/.style={font=\tiny, text=black!65},
    arrow/.style={-Latex, line width=0.95pt, draw=black!65},
    darrow/.style={-Latex, line width=0.85pt, draw=black!42, dashed},
    imtag/.style={font=\tiny\bfseries, text=white, inner sep=1.2pt, rounded corners=1pt, anchor=north west},
    tbox/.style={draw=truthcol, line width=1.1pt, minimum width=1.05cm, minimum height=0.62cm, align=center, font=\small},
    pbox/.style={draw=predcol, line width=1.1pt, minimum width=1.05cm, minimum height=0.62cm, align=center, font=\small},
    eslot/.style={draw=black!35, dashed, line width=0.9pt, minimum width=1.05cm, minimum height=0.62cm, align=center, font=\small, text=black!45},
    dot/.style={circle, inner sep=0pt, minimum size=7pt}
]

% ---- Layout -------------------------------------------------------------
\def\gap{0.26}
\def\wA{6.20}
\def\wB{4.95}
\def\wC{5.15}
\def\wD{9.50}   % bottom-left: confusion matrix (shrunk)
\def\wE{7.06}   % bottom-right: where are the TN?  (\wD + gap + \wE = top-row width)
\def\hTop{5.10}
\def\hBot{4.15}

\coordinate (NW) at (0,0);
\node[panel, minimum width=\wA cm, minimum height=\hTop cm, anchor=north west] (P1) at (NW) {};
\node[panel, minimum width=\wB cm, minimum height=\hTop cm, anchor=north west] (P2) at ([xshift=\wA cm + \gap cm]NW) {};
\node[panel, minimum width=\wC cm, minimum height=\hTop cm, anchor=north west] (P3) at ([xshift=\wA cm + \wB cm + 2*\gap cm]NW) {};
\node[panel, minimum width=\wD cm, minimum height=\hBot cm, anchor=north west] (P4) at ([yshift=-\hTop cm - \gap cm]NW) {};
\node[panel, minimum width=\wE cm, minimum height=\hBot cm, anchor=north west] (P5) at ([xshift=\wD cm + \gap cm, yshift=-\hTop cm - \gap cm]NW) {};

\node[ptitle] at ([xshift=0.18cm,yshift=-0.16cm]P1.north west) {1. Inputs};
\node[ptitle] at ([xshift=0.18cm,yshift=-0.16cm]P2.north west) {2. Assignment};
\node[ptitle] at ([xshift=0.18cm,yshift=-0.16cm]P3.north west) {3. Classification view};
\node[ptitle] at ([xshift=0.18cm,yshift=-0.16cm]P4.north west) {4. Confusion matrix};
\node[ptitle] at ([xshift=0.18cm,yshift=-0.16cm]P5.north west) {5. Where are the \TN?};

% =====================  Panel 1: image, truth, predictions  =============
\node[body, anchor=south west]  at ([xshift=0.70cm,yshift=-1.20cm]P1.north west) {Image};
\node[body, text=truthcol, anchor=south] at ([xshift=3.10cm,yshift=-1.20cm]P1.north west) {Ground truth};
\node[body, text=predcol, anchor=south] at ([xshift=5.02cm,yshift=-1.20cm]P1.north west) {Predictions};

\node[anchor=north west, inner sep=0pt] (I1) at ([xshift=0.24cm,yshift=-1.14cm]P1.north west)
    {\includegraphics[width=\Wd cm]{cat-and-dog.jpg}};
\node[anchor=north west, inner sep=0pt] (I2) at ([xshift=2.18cm,yshift=-1.14cm]P1.north west)
    {\includegraphics[width=\Wd cm]{cat-and-dog.jpg}};
\node[anchor=north west, inner sep=0pt] (I3) at ([xshift=4.12cm,yshift=-1.14cm]P1.north west)
    {\includegraphics[width=\Wd cm]{cat-and-dog.jpg}};

% Ground-truth boxes: T1 = dog, T2 = cat
\draw[truthcol, line width=1.2pt] ([xshift=0.054cm,yshift=-0.405cm]I2.north west) rectangle ([xshift=1.098cm,yshift=-2.592cm]I2.north west);
\draw[truthcol, line width=1.2pt] ([xshift=0.957cm,yshift=-1.188cm]I2.north west) rectangle ([xshift=1.741cm,yshift=-2.511cm]I2.north west);
\node[imtag, fill=truthcol] at ([xshift=0.054cm,yshift=-0.405cm]I2.north west) {T1};
\node[imtag, fill=truthcol] at ([xshift=0.887cm,yshift=-1.188cm]I2.north west) {T2};

% Predicted boxes: P1 = good dog detection, P2 = spurious box on grass
\draw[predcol, line width=1.2pt] ([xshift=0.054cm,yshift=-0.459cm]I3.north west) rectangle ([xshift=1.053cm,yshift=-2.565cm]I3.north west);
\draw[predcol, line width=1.2pt] ([xshift=1.260cm,yshift=-0.162cm]I3.north west) rectangle ([xshift=1.656cm,yshift=-0.810cm]I3.north west);
\node[imtag, fill=predcol] at ([xshift=0.054cm,yshift=-0.459cm]I3.north west) {P1};
\node[imtag, fill=predcol, anchor=north east] at ([xshift=1.656cm,yshift=-0.162cm]I3.north west) {P2};

% Caption + image credit, kept inside Panel 1's width so it doesn't bleed across.
\node[note, align=left, text width=5.70cm, anchor=south west] at ([xshift=0.26cm,yshift=0.52cm]P1.south west) {%
Dog found (T1/P1), grass falsely flagged (P2), cat missed (T2).};
\node[font=\tiny, text=black!42, align=left, text width=5.70cm, anchor=south west] at ([xshift=0.26cm,yshift=0.14cm]P1.south west) {%
Photo: publicdomainpictures.net \#24076, CC0.};

% =====================  Panel 2: assignment  ===========================
\node[body, align=center, anchor=north] at ([yshift=-0.52cm]P2.north)
    {Match each prediction to a truth box};
\node[note, anchor=north] at ([yshift=-0.86cm]P2.north)
    {e.g.\ one-to-one by IoU overlap};

\coordinate (Q) at (P2.north west);
\node[body, text=truthcol, anchor=south] at ($(Q)+(1.15,-1.76)$) {Truth};
\node[body, text=predcol,  anchor=south] at ($(Q)+(3.05,-1.76)$) {Prediction};
\node[body, anchor=south] at ($(Q)+(4.45,-1.76)$) {Result};

% Row 1: matched -> TP
\node[tbox] (T1) at ($(Q)+(1.15,-2.08)$) {$T_1$};
\node[pbox] (P1b) at ($(Q)+(3.05,-2.08)$) {$P_1$};
\draw[arrow] (T1.east) -- (P1b.west);
\node at ($(Q)+(4.45,-2.08)$) {\TP};

% Row 2: truth with no prediction -> FN
\node[tbox] (T2) at ($(Q)+(1.15,-3.02)$) {$T_2$};
\node[eslot] (E1) at ($(Q)+(3.05,-3.02)$) {$\varnothing$};
\draw[darrow] (T2.east) -- (E1.west);
\node at ($(Q)+(4.45,-3.02)$) {\FN};

% Row 3: prediction with no truth -> FP
\node[eslot] (E2) at ($(Q)+(1.15,-3.96)$) {$\varnothing$};
\node[pbox] (P2b) at ($(Q)+(3.05,-3.96)$) {$P_2$};
\draw[darrow] (P2b.west) -- (E2.east);
\node at ($(Q)+(4.45,-3.96)$) {\FP};

\node[note, anchor=south] at ([yshift=0.10cm]P2.south)
    {solid = matched\quad\textbullet\quad dashed = no match ($\varnothing$)};

% =====================  Panel 3: classification view (outcome coloring)  ====
% Same truth/prediction pair as Panel 1, but each box is now recolored by the
% outcome it earned: matched dog = TP (green), missed cat = FN (orange),
% spurious grass box = FP (red).  No TN appears -- which is the whole point.
\def\Ws{1.60}   % small copy width  (cm)
\def\Hs{2.40}   % small copy height (= 1.60 * 1.5)

\node[note, text=truthcol, anchor=south] at ([xshift=1.36cm,yshift=-1.15cm]P3.north west) {Ground truth};
\node[note, text=predcol, anchor=south] at ([xshift=3.74cm,yshift=-1.15cm]P3.north west) {Predictions};

\node[anchor=north west, inner sep=0pt] (J1) at ([xshift=0.56cm,yshift=-1.10cm]P3.north west)
    {\includegraphics[width=\Ws cm]{cat-and-dog.jpg}};
\node[anchor=north west, inner sep=0pt] (J2) at ([xshift=2.94cm,yshift=-1.10cm]P3.north west)
    {\includegraphics[width=\Ws cm]{cat-and-dog.jpg}};

% Truth image: dog -> TP (green), cat -> FN (orange)
\draw[kwgreen,  line width=1.2pt] ([xshift=0.048cm,yshift=-0.360cm]J1.north west) rectangle ([xshift=0.976cm,yshift=-2.304cm]J1.north west);
\draw[kworange, line width=1.2pt] ([xshift=0.851cm,yshift=-1.056cm]J1.north west) rectangle ([xshift=1.548cm,yshift=-2.232cm]J1.north west);
\node[imtag, fill=kwgreen]  at ([xshift=0.048cm,yshift=-0.360cm]J1.north west) {TP};
\node[imtag, fill=kworange] at ([xshift=0.851cm,yshift=-1.056cm]J1.north west) {FN};

% Prediction image: dog -> TP (green), grass box -> FP (red)
\draw[kwgreen, line width=1.2pt] ([xshift=0.048cm,yshift=-0.408cm]J2.north west) rectangle ([xshift=0.936cm,yshift=-2.280cm]J2.north west);
\draw[kwred,   line width=1.2pt] ([xshift=1.120cm,yshift=-0.144cm]J2.north west) rectangle ([xshift=1.472cm,yshift=-0.720cm]J2.north west);
\node[imtag, fill=kwgreen]                  at ([xshift=0.048cm,yshift=-0.408cm]J2.north west) {TP};
\node[imtag, fill=kwred, anchor=north east] at ([xshift=1.472cm,yshift=-0.144cm]J2.north west) {FP};

% Legend: define each outcome and its color.
\node[note, anchor=north west] at ([xshift=0.34cm,yshift=-3.62cm]P3.north west) {Each match or leftover is one of:};
\node[body, align=left, text width=2.35cm, anchor=north west] at ([xshift=0.30cm,yshift=-3.92cm]P3.north west) {%
\tikz{\node[dot,fill=kwgreen]{};}\,\TP\ matched pair\\[0.55em]
\tikz{\node[dot,fill=kwred]{};}\,\FP\ false box};
\node[body, align=left, text width=2.35cm, anchor=north west] at ([xshift=2.75cm,yshift=-3.92cm]P3.north west) {%
\tikz{\node[dot,fill=kworange]{};}\,\FN\ missed truth\\[0.55em]
\tikz{\node[dot,fill=kwblue]{};}\,\TN\ empty region};

% =====================  Panel 4: confusion matrix  =====================
\coordinate (CMNW) at ([xshift=1.20cm,yshift=-1.18cm]P4.north west);
\draw[fill=kwgreen!20,  draw=white] (CMNW) rectangle ++(1.98,-1.22);
\draw[fill=kworange!22, draw=white] ([xshift=2.12cm]CMNW) rectangle ++(1.98,-1.22);
\draw[fill=kwred!16,    draw=white] ([yshift=-1.34cm]CMNW) rectangle ++(1.98,-1.22);
\draw[fill=kwblue!16,   draw=white] ([xshift=2.12cm,yshift=-1.34cm]CMNW) rectangle ++(1.98,-1.22);

\node[body] at ($(CMNW)+(0.99,-0.40)$) {\TP};
\node[note,align=center] at ($(CMNW)+(0.99,-0.84)$) {matched\\(here: 1)};
\node[body] at ($(CMNW)+(3.11,-0.40)$) {\FN};
\node[note,align=center] at ($(CMNW)+(3.11,-0.84)$) {missed truth\\(here: 1)};
\node[body] at ($(CMNW)+(0.99,-1.73)$) {\FP};
\node[note,align=center] at ($(CMNW)+(0.99,-2.17)$) {spurious box\\(here: 1)};
\node[body] at ($(CMNW)+(3.11,-1.73)$) {\TN};
\node[note,align=center] at ($(CMNW)+(3.11,-2.17)$) {correctly empty\\(here: 0)};

\node[body, anchor=south] at ($(CMNW)+(2.05,0.10)$) {Predicted};
\node[note, anchor=south] at ($(CMNW)+(0.99,-0.08)$) {positive};
\node[note, anchor=south] at ($(CMNW)+(3.11,-0.08)$) {negative};
\node[body, rotate=90] at ($(CMNW)+(-0.78,-1.28)$) {Real};
\node[note, rotate=90, anchor=south] at ($(CMNW)+(-0.16,-0.61)$) {positive};
\node[note, rotate=90, anchor=south] at ($(CMNW)+(-0.16,-1.95)$) {negative};

% --- Keep the "view detection as classification" takeaway beside the matrix.
\node[body, align=left, text width=3.45cm, anchor=north west] at ([xshift=5.70cm,yshift=-1.18cm]P4.north west) {
%\textbf{View detection as classification.}\\[0.3em]
With the assignment fixed, each box is one cell of this matrix: a matched pair is a \TP, a prediction with no truth a \FP, and a truth with no prediction a \FN.\\[0.3em]
Our small example gives one of each. But notice that \TN{} is zero and typically uncounted.};

% =====================  Panel 5: where are the TN?  ====================
% One more copy of the image.  The two real objects (T1, T2) are boxed; every
% OTHER box you could draw is a true negative.  The TN grid is masked so it
% never overlaps T1 or T2 -- the empty background is what "explodes".
\def\Wt{2.00}   % TN image width  (cm)
\def\Ht{3.00}   % TN image height (= 2.00 * 1.5)
\node[anchor=north west, inner sep=0pt] (IT) at ([xshift=0.40cm,yshift=-1.00cm]P5.north west)
    {\includegraphics[width=\Wt cm]{cat-and-dog.jpg}};

% Dense grid of candidate background boxes, skipping any cell that touches an object.
\foreach \i in {0,...,7}{
  \foreach \j in {0,...,11}{
    \pgfmathsetmacro{\nxa}{\i/8}\pgfmathsetmacro{\nxb}{(\i+1)/8}
    \pgfmathsetmacro{\nya}{\j/12}\pgfmathsetmacro{\nyb}{(\j+1)/12}
    \pgfmathtruncatemacro{\hit}{ ((\nxa<0.610)&&(\nxb>0.030)&&(\nya<0.960)&&(\nyb>0.150)) || ((\nxa<0.967)&&(\nxb>0.532)&&(\nya<0.930)&&(\nyb>0.440)) ? 1 : 0 }
    \ifnum\hit=0
      \pgfmathsetmacro{\bxa}{(\i+0.16)/8*\Wt}\pgfmathsetmacro{\bxb}{(\i+0.84)/8*\Wt}
      \pgfmathsetmacro{\bya}{-(\j+0.16)/12*\Ht}\pgfmathsetmacro{\byb}{-(\j+0.84)/12*\Ht}
      \draw[kwblue, line width=0.25pt, fill=kwblue, fill opacity=0.13]
        ([xshift=\bxa cm,yshift=\bya cm]IT.north west) rectangle ([xshift=\bxb cm,yshift=\byb cm]IT.north west);
    \fi
  }}
% A couple of larger background boxes, to show that ANY size counts too.
\draw[kwblue, line width=0.5pt, fill=kwblue, fill opacity=0.07] ([xshift=0.10cm,yshift=-0.06cm]IT.north west) rectangle ([xshift=1.10cm,yshift=-0.39cm]IT.north west);
\draw[kwblue, line width=0.5pt, fill=kwblue, fill opacity=0.07] ([xshift=1.28cm,yshift=-0.15cm]IT.north west) rectangle ([xshift=1.94cm,yshift=-1.20cm]IT.north west);

% The two real objects, for reference.
\draw[truthcol, line width=1.1pt] ([xshift=0.06cm,yshift=-0.45cm]IT.north west) rectangle ([xshift=1.22cm,yshift=-2.88cm]IT.north west);
\draw[truthcol, line width=1.1pt] ([xshift=1.064cm,yshift=-1.32cm]IT.north west) rectangle ([xshift=1.934cm,yshift=-2.79cm]IT.north west);
\node[imtag, fill=truthcol] at ([xshift=0.06cm,yshift=-0.45cm]IT.north west) {T1};
\node[imtag, fill=truthcol] at ([xshift=1.064cm,yshift=-1.32cm]IT.north west) {T2};
\node[imtag, fill=kwblue] at ([xshift=0.04cm,yshift=-0.04cm]IT.north west) {TN};

\node[body, align=left, text width=4.05cm, anchor=north west] at ([xshift=2.75cm,yshift=-1.05cm]P5.north west) {
Boxing the two objects leaves \emph{everything else}: each box you could draw that matches no real object is a \TN --- at every position and \emph{every} size.\\[0.35em]
There is no natural limit, so the usual metrics ignore \TN. This paper instead asks what the score approaches as \TN${}\to\infty$.};

\end{tikzpicture}
\end{document}
